\section{AHSME 1950}
        \begin{enumerate}
        	\item (\textit{AHSME 1950}) If $64$ is divided into three parts proportional to $2$, $4$, and $6$, the smallest part is:


 $\textbf{(A)}\ 5\frac{1}{3}\qquad\textbf{(B)}\ 11\qquad\textbf{(C)}\ 10\frac{2}{3}\qquad\textbf{(D)}\ 5\qquad\textbf{(E)}\ \text{None of these answers}$

 \textbf{Answer: }C



\textbf{Solution 1}

Given,The ratios are 2:4:6 The number is 64

 So, the sum of the ratios is 12 

 So, the first number is  
 \[=\frac{64*2}{12}\]=
 \[\frac{32}{3}\] So, the second number is 
 \[=\frac{64*4}{12}\]=
 \[\frac{64}{3}\] 

 So, the third number is
 \[=\frac{64*6}{12}\]=32

 \begin{verbatim}
We can see that the smallest is
 \[=\frac{32}{3}\]
\end{verbatim}


 \[x=\frac{32}{3}=10 \frac{2}{3}\]which is $\boxed{\textbf{(C)}}$.

 



\textbf{Solution 2}

If the three numbers are in proportion to $2:4:6$, then they should also be in proportion to $1:2:3$. This implies that the three numbers can be expressed as $x$, $2x$, and $3x$. Add these values together to get:
 \[x+2x+3x=6x=64\]Divide each side by 6 and get that 
 \[x=\frac{64}{6}=\frac{32}{3}=10 \frac{2}{3}\]which is $\boxed{\textbf{(C)}}$.

 


	\item (\textit{AHSME 1950}) Let $R=gS-4$. When $S=8$, $R=16$. When $S=10$, $R$ is equal to:


 $\textbf{(A)}\ 11\qquad\textbf{(B)}\ 14\qquad\textbf{(C)}\ 20\qquad\textbf{(D)}\ 21\qquad\textbf{(E)}\ \text{None of these}$

 \textbf{Answer: }D



\textbf{Solution}

Our first procedure is to find the value of $g$. With the given variables' values, we can see that $8g-4=16$ so $g=\frac{20}{8}=\frac{5}{2}$.

 With that, we can replace $g$ with $\frac{5}{2}$. When $S=10$, we can see that $10\times\frac{5}{2}-4=\frac{50}{2}-4=25-4=\boxed{\text{(D) 21}}$.

 


	\item (\textit{AHSME 1950}) The sum of the roots of the equation $4x^{2}+5-8x=0$ is equal to:


 $\textbf{(A)}\ 8\qquad\textbf{(B)}\ -5\qquad\textbf{(C)}\ -\frac{5}{4}\qquad\textbf{(D)}\ -2\qquad\textbf{(E)}\ \text{None of these}$

 \textbf{Answer: }E



\textbf{Solution}

We can divide by 4 to get:$x^{2}-2x+\dfrac{5}{4}=0.$

 The Vieta's formula states that in quadratic equation $ax^2+bx+c$, the sum of the roots of the equation is $-\frac{b}{a}$.Using Vieta's formula, we find that the roots add to $2$ or $\boxed{\textbf{(E)}\ \text{None of these}}$.

 



\textbf{Solution 2}

We can divide by 4 to get:$x^{2}-2x+\dfrac{5}{4}=0.$

 By using quadratic formula you can find the roots are \(x=\frac{8\pm \sqrt{-16}}{8}\) adding the roots together to get 2.

 or $\boxed{\textbf{(E)}\ \text{None of these}}$.

 


	\item (\textit{AHSME 1950}) Reduced to lowest terms, $\frac{a^{2}-b^{2}}{ab} - \frac{ab-b^{2}}{ab-a^{2}}$ is equal to:


 $\textbf{(A)}\ \frac{a}{b}\qquad\textbf{(B)}\ \frac{a^{2}-2b^{2}}{ab}\qquad\textbf{(C)}\ a^{2}\qquad\textbf{(D)}\ a-2b\qquad\textbf{(E)}\ \text{None of these}$

 \textbf{Answer: }A



\textbf{Solution}

We start off by factoring the second fraction.

 
 \[-\frac{ab-b^2}{ab-a^2} = -\frac{b(a-b)}{-a(a-b)} = \frac{b}{a}.\]

 Now create a common denominator and simplify.

 
 \[\frac{a^2-b^2}{ab}+\frac{b}{a}=\frac{a^2-b^2}{ab}+\frac{b^2}{ab} = \frac{a^2}{ab} = \boxed{\mathrm{(A) }\frac{a}{b}}\]

 obs: Assume that $a \not = 0, b\not = 0, a\not = b$.

 


	\item (\textit{AHSME 1950}) If five geometric means are inserted between $8$ and $5832$, the fifth term in the geometric series is:


 $\textbf{(A)}\ 648\qquad\textbf{(B)}\ 832\qquad\textbf{(C)}\ 1168\qquad\textbf{(D)}\ 1944\qquad\textbf{(E)}\ \text{None of these}$

 \textbf{Answer: }A



\textbf{Solution}

We can let the common ratio of the geometric sequence be $r$. $5832$ is given to be the seventh term in the geometric sequence as there are five terms between it and $a_1$ if we consider $a_1=8$.By the formula for each term in a geometric sequence, we find that $a_n=a_1r^{n-1}$ or $(5382)=(8)r^6$We divide by eight to find:$r^6=729$$r=\pm 3$

 Because $a_2$  will not be between $8$ and $5832$ if $r=-3$ we can discard it as an extraneous solution. We find $r=3$ and $a_5=a_1r^4=(8)(3)^4=\boxed{\textbf{(A)}\ 648}$



	\item (\textit{AHSME 1950}) The values of $y$ which will satisfy the equations 
 \[\begin{array}{rcl} 2x^{2}+6x+5y+1&=&0\\ 2x+y+3&=&0 \end{array}\] may be found by solving:


 $\textbf{(A)}\ y^{2}+14y-7=0\qquad\textbf{(B)}\ y^{2}+8y+1=0\qquad\textbf{(C)}\ y^{2}+10y-7=0\qquad\\ \textbf{(D)}\ y^{2}+y-12=0\qquad \textbf{(E)}\ \text{None of these equations}$

 \textbf{Answer: }C



\textbf{Solution}

If we solve the second equation for $x$ in terms of $y$, we find $x=-\dfrac{y+3}{2}$ which we can substitute to find:

 
 \[2(-\dfrac{y+3}{2})^2+6(-\dfrac{y+3}{2})+5y+1=0\]

 Multiplying by two and simplifying, we find:

 
 \begin{align*} 2\cdot[2(-\dfrac{y+3}{2})^2+6(-\dfrac{y+3}{2})+5y+1]&=2\cdot 0\\ (y+3)^2 -6y-18+10y+2&=0\\ y^2+6y+9-6y-18+10y+2&=0\\ y^2+10y-7&=0 \end{align*}


 Therefore the answer is $\boxed{\textbf{(C)}\ y^{2}+10y-7=0}$



	\item (\textit{AHSME 1950}) If the digit $1$ is placed after a two digit number whose tens' digit is $t$, and units' digit is $u$, the new number is:


 $\textbf{(A)}\ 10t+u+1\qquad\textbf{(B)}\ 100t+10u+1\qquad\textbf{(C)}\ 1000t+10u+1\qquad\textbf{(D)}\ t+u+1\qquad\\ \textbf{(E)}\ \text{None of these answers}$

 \textbf{Answer: }B



\textbf{Solution}

By placing the digit $1$ after a two digit number, you are changing the units place to $1$ and moving everything else up a place. Therefore the answer is $\boxed{\textbf{(B)}\ 100t+10u+1}$.

 


	\item (\textit{AHSME 1950}) If the radius of a circle is increased $100\%$, the area is increased:


 $\textbf{(A)}\ 100\%\qquad\textbf{(B)}\ 200\%\qquad\textbf{(C)}\ 300\%\qquad\textbf{(D)}\ 400\%\qquad\textbf{(E)}\ \text{By none of these}$

 \textbf{Answer: }C



\textbf{Solution}

Increasing by $100\%$ is the same as doubling the radius. If we let $r$ be the radius of the old circle, then the radius of the new circle is $2r.$

 Since the area of the circle is given by the formula $\pi r^2,$ the area of the new circle is $\pi (2r)^2 = 4\pi r^2.$ The area is quadrupled, or increased by $\boxed{\mathrm{(C) }300\%.}$



	\item (\textit{AHSME 1950}) The largest area of a triangle that can be inscribed in a semi-circle whose radius is $r$ is:


 $\textbf{(A)}\ r^{2}\qquad\textbf{(B)}\ r^{3}\qquad\textbf{(C)}\ 2r^{2}\qquad\textbf{(D)}\ 2r^{3}\qquad\textbf{(E)}\ \frac{1}{2}r^{2}$

 \textbf{Answer: }A



\textbf{Solution}

The area of a triangle is $\frac12 bh.$ To maximize the base, let it be equal to the diameter of the semi circle, which is equal to $2r.$ To maximize the height, or altitude, choose the point directly in the middle of the arc connecting the endpoints of the diameter. It is equal to $r.$ Therefore the area is $\frac12 \cdot 2r \cdot r = \boxed{\mathrm{(A) }r^2}$.

 


	\item (\textit{AHSME 1950}) After rationalizing the numerator of $\frac{\sqrt{3}-\sqrt{2}}{\sqrt{3}}$, the denominator in simplest form is:


 $\textbf{(A)}\ \sqrt{3}(\sqrt{3}+\sqrt{2})\qquad\textbf{(B)}\ \sqrt{3}(\sqrt{3}-\sqrt{2})\qquad\textbf{(C)}\ 3-\sqrt{3}\sqrt{2}\qquad\\ \textbf{(D)}\ 3+\sqrt6\qquad\textbf{(E)}\ \text{None of these answers}$

 \textbf{Answer: }D



\textbf{Solution}

To rationalize the numerator, multiply by the conjugate.

 
 \[\frac{\sqrt3 - \sqrt2}{\sqrt3}\cdot \frac{\sqrt3 + \sqrt2}{\sqrt3 + \sqrt2} = \frac{1}{3+\sqrt6}.\]

 The denominator is $\boxed{\mathrm{(D)}\text{ } 3+\sqrt6 }.$



	\item (\textit{AHSME 1950}) If in the formula $C =\frac{en}{R+nr}$, $n$ is increased while $e$, $R$ and $r$ are kept constant, then $C$:


 $\textbf{(A)}\ \text{decreases}\qquad\textbf{(B)}\ \text{increases}\qquad\textbf{(C)}\ \text{remains constant}\qquad\textbf{(D)}\ \text{increases and then decreases}\qquad\\ \textbf{(E)}\ \text{decreases and then increases}$

 \textbf{Answer: }B



\textbf{Solution}

Divide both the numerator and denominator by $n$, to get $C=\frac{e}{\frac{R}{n}+r}$.  If $n$ increases then the denominator decreases; so that $C$ $\boxed{\mathbf{(B)}\text{ }\mathrm{ increases}.}$



	\item (\textit{AHSME 1950}) As the number of sides of a polygon increases from $3$ to $n$, the sum of the exterior angles formed by extending each side in succession:


 $\textbf{(A)}\ \text{increases}\qquad\textbf{(B)}\ \text{decreases}\qquad\textbf{(C)}\ \text{remains constant}\qquad\textbf{(D)}\ \text{cannot be predicted}\qquad\\ \textbf{(E)}\ \text{becomes }(n-3)\text{ straight angles}$

 \textbf{Answer: }C



\textbf{Solution}

The angles of the exterior polygon always add up to 180 degrees, therefore, the answer remains constant. So (C)

 


 First of all, the sum of exterior angles of a polygon is always 360°.Second of all, in MY opinion, the question is kind of arbitrary and I don't know which option you're talking about. Someone probably has cut the question.                                                                   ~Randommathdoerfromindia

 He is using increasing angles. To do what?

 


	\item (\textit{AHSME 1950}) The roots of $(x^{2}-3x+2)(x)(x-4)=0$ are:


 $\textbf{(A)}\ 4\qquad\textbf{(B)}\ 0\text{ and }4\qquad\textbf{(C)}\ 1\text{ and }2\qquad\textbf{(D)}\ 0,1,2\text{ and }4\qquad\textbf{(E)}\ 1,2\text{ and }4$

 \textbf{Answer: }D



\textbf{Solution}

Factor $x^2-3x+2$ to get $(x-2)(x-1).$ The roots are $\boxed{\mathrm{(D)}\ 0,1,2\text{ and }4.}$



	\item (\textit{AHSME 1950}) For the simultaneous equations 
 \[2x-3y=8\]
 \[6y-4x=9\]
$\textbf{(A)}\ x=4,y=0\qquad\textbf{(B)}\ x=0,y=\frac{3}{2}\qquad\textbf{(C)}\ x=0,y=0\qquad\\ \textbf{(D)}\ \text{There is no solution}\qquad\textbf{(E)}\ \text{There are an infinite number of solutions}$

 \textbf{Answer: }D



\textbf{Solution}

Try to solve this system of equations using the elimination method.

 
 \begin{align*} -2(2x-3y)&=-2(8)\\ 6y-4x&=-16\\ 6y-4x&=9\\ 0&=-7 \end{align*}


 Something is clearly contradictory so $\boxed{\mathrm{(D)}\text{ There is no solution}.}$

 Alternatively, note that the second equation is a multiple of the first except that the constants don't match up. So, there is no solution.

 


	\item (\textit{AHSME 1950}) The real factors of $x^2+4$ are:


 $\textbf{(A)}\ (x^{2}+2)(x^{2}+2)\qquad\textbf{(B)}\ (x^{2}+2)(x^{2}-2)\qquad\textbf{(C)}\ x^{2}(x^{2}+4)\qquad\\ \textbf{(D)}\ (x^{2}-2x+2)(x^{2}+2x+2)\qquad\textbf{(E)}\ \text{Non-existent}$

 \textbf{Answer: }E



\textbf{Solution 1}

Notice that $x^2 + 4 = x^2 - (-4) = x^2 - (2i)^2 = (x+2i)(x-2i)$. Therefore, $x^2-4$ cannot be further factored into polynomials with real coefficients, and the answer is $\boxed{\mathrm{(E)}\text{ Non-existent.}}$




\textbf{Solution 2 (Answer Choices)}

Notice that the first four answer choices each have two factors of degree $2$.  Therefore, they each have degree $4$ and are incorrect.  Thus $\boxed{\mathbf{(E)}\text{ Non-existent}}$ is the correct answer.

 



\textbf{Solution 3}

Let's try to find all real roots of $x^2+4=0$. If there was a real root (call it $r$), then it would satisfy $r^2+4=0$. Rearranging gives that $r^2=-4$. Therefore a real root $r$ of $x^2+4=0$ would satisfy $r^2<0$. However, the \href{https://artofproblemsolving.com/wiki/index.php?title=Trivial\_Inequality}{Trivial Inequality} states that no real $r$ satisfies $r^2<0$, which must mean that there are no real roots $r$; they are $\boxed{\mathrm{(E)}\text{ Non-existent.}}$




\textbf{Solution 4 (Easier Solution)}

Let us try to factor $x^2+4.$ We can see that the first term has variables only and the second has a constant only. Therefore, we cannot factor out anything and our answer is $\boxed{\mathrm{(E)}\text{ Non-existent.}}$



	\item (\textit{AHSME 1950}) The number of terms in the expansion of $[(a+3b)^{2}(a-3b)^{2}]^{2}$ when simplified is:


 $\textbf{(A)}\ 4\qquad\textbf{(B)}\ 5\qquad\textbf{(C)}\ 6\qquad\textbf{(D)}\ 7\qquad\textbf{(E)}\ 8$

 \textbf{Answer: }B



\textbf{Solution}

Use properties of exponents to move the squares outside the brackets use difference of squares.

 
 \[[(a+3b)(a-3b)]^4 = (a^2-9b^2)^4\]

 Using the binomial theorem, we can see that the number of terms is $\boxed{\mathrm{(B)}\ 5}$.

 


	\item (\textit{AHSME 1950}) The formula which expresses the relationship between $x$ and $y$ as shown in the accompanying table is:


 
 \[\begin{tabular}[t]{|c|c|c|c|c|c|}\hline x&0&1&2&3&4\\\hline y&100&90&70&40&0\\\hline\end{tabular}\]
$\textbf{(A)}\ y=100-10x\qquad\textbf{(B)}\ y=100-5x^{2}\qquad\textbf{(C)}\ y=100-5x-5x^{2}\qquad\\ \textbf{(D)}\ y=20-x-x^{2}\qquad\textbf{(E)}\ \text{None of these}$

 \textbf{Answer: }C



\textbf{Solution}

Plug in the points $(0,100)$ and $(4,0)$ into each equation. The only one that works for both points is $\mathrm{(C)}.$ Plug in the rest of the points to confirm the answer is indeed $\boxed{\mathrm{(C)}\ y=100-5x-5x^2}$.

 


	\item (\textit{AHSME 1950}) Of the following 


 (1) $a(x-y)=ax-ay$


 (2) $a^{x-y}=a^x-a^y$


 (3) $\log (x-y)=\log x-\log y$


 (4) $\frac{\log x}{\log y}=\log{x}-\log{y}$


 (5) $a(xy)=ax \cdot ay$
$\textbf{(A)}\ \text{Only 1 and 4 are true}\qquad\\\textbf{(B)}\ \text{Only 1 and 5 are true}\qquad\\\textbf{(C)}\ \text{Only 1 and 3 are true}\qquad\\\textbf{(D)}\ \text{Only 1 and 2 are true}\qquad\\\textbf{(E)}\ \text{Only 1 is true}$

 \textbf{Answer: }E



\textbf{Solution}

The distributive property doesn't apply to logarithms or in the ways illustrated, and only applies to addition and subtraction. Also, $a^{x-y} = \frac{a^x}{a^y}$, so $\boxed{\textbf{(E)} \text{ Only 1 is true}}$.

 


	\item (\textit{AHSME 1950}) If $m$ men can do a job in $d$ days, then $m+r$ men can do the job in:


 $\textbf{(A)}\ d+r\text{ days}\qquad\textbf{(B)}\ d-r\text{ days}\qquad\textbf{(C)}\ \frac{md}{m+r}\text{ days}\qquad\textbf{(D)}\ \frac{d}{m+r}\text{ days}\qquad\textbf{(E)}\ \text{None of these}$

 \textbf{Answer: }C



\textbf{Solution}

The number of men is inversely proportional to the number of days the job takes. Thus, if $m$ men can do a job in $d$ days, we have that it will take $md$ days for $1$ man to do the job. Thus, $m + r$ men can do the job in $\boxed{\textbf{(C)}\ \frac{md}{m+r}\text{ days}}$.

 


	\item (\textit{AHSME 1950}) When $x^{13}+1$ is divided by $x-1$, the remainder is:


 $\textbf{(A)}\ 1\qquad\textbf{(B)}\ -1\qquad\textbf{(C)}\ 0\qquad\textbf{(D)}\ 2\qquad\textbf{(E)}\ \text{None of these answers}$

 \textbf{Answer: }D



\textbf{Solution 1}

Using synthetic division, we get that the remainder is $\boxed{\textbf{(D)}\ 2}$.

 



\textbf{Solution 2}

By the remainder theorem, the remainder is equal to the expression $x^{13}+1$ when $x=1.$ This gives the answer of $\boxed{(\mathrm{D})\ 2.}$




\textbf{Solution 3}

Note that $x^{13} - 1 = (x - 1)(x^{12} + x^{11} \cdots + 1)$, so $x^{13} - 1$ is divisible by $x-1$, meaning $(x^{13} - 1) + 2$ leaves a remainder of $\boxed{\mathrm{(D)}\ 2.}$




\textbf{Solution 4}

Say the quotient of the polynomial is \( q \) and the remainder is \( r \). Then we have

 
 \[x^{13} + 1 = (x - 1)q + r\]

 To eliminate \( q \), we put \( x = 1 \) in which we get

 
 \[1^{13} + 1 = r\]

 \[1 + 1 = r\]

 \[2 = r\]

 Thus the answer is $\boxed{(\mathrm{D})\ 2.}$



	\item (\textit{AHSME 1950}) The volume of a rectangular solid each of whose side, front, and bottom faces are $12\text{ in}^{2}$, $8\text{ in}^{2}$, and $6\text{ in}^{2}$ respectively is:


 $\textbf{(A)}\ 576\text{ in}^{3}\qquad\textbf{(B)}\ 24\text{ in}^{3}\qquad\textbf{(C)}\ 9\text{ in}^{3}\qquad\textbf{(D)}\ 104\text{ in}^{3}\qquad\textbf{(E)}\ \text{None of these}$

 \textbf{Answer: }B



\textbf{Solution}

If the sidelengths of the cubes are expressed as $a, b,$ and $c,$ then we can write three equations:

 
 \[ab=12, bc=8, ac=6.\]

 The volume is $abc.$ Notice symmetry in the equations. We can find $abc$ my multiplying all the equations and taking the positive square root.

 
 \begin{align*} (ab)(bc)(ac) &= (12)(8)(6)\\ a^2b^2c^2 &= 576\\ abc &= \boxed{\mathrm{(B)}\ 24} \end{align*}




	\item (\textit{AHSME 1950}) Successive discounts of $10\%$ and $20\%$ are equivalent to a single discount of:


 $\textbf{(A)}\ 30\%\qquad\textbf{(B)}\ 15\%\qquad\textbf{(C)}\ 72\%\qquad\textbf{(D)}\ 28\%\qquad\textbf{(E)}\ \text{None of these}$

 \textbf{Answer: }D



\textbf{Solution 1 (Kind of Lame)}

Without loss of generality, assume something costs $100$ dollars. Then with each successive discount, it would cost $90$ dollars, then $72$ dollars. This amounts to a total of $28$ dollars off, so the single discount would be $\boxed{\mathrm{(D)}\ 28\%.}$




\textbf{Solution 2 (Quite similar to sollution 1)}

Total discount if the Product is $100$ dollars  is  10/100 of 100 +  20/100 of (100 - 10/100 of 100)

 \begin{verbatim}
												= 10/100  * 100 +  20/100  * (100 - 10/100  * 100)                                                  = 10 + 2/10 * 90                                                 = 10 + 9*2                                                 = 10+18                                                 = 28
\end{verbatim}

So discount is $28$ dollarsso the single discount would be $\boxed{\mathrm{(D)}\ 28\%.}$




\textbf{Solution 3 (Technical)}

Let the object cost $x$ dollars. After the $10\%$ discount, it's worth $(1-10\%)x=0.9x$ dollars. After a $20\%$ discount on that, it's worth $(1-20\%)(0.9x)=0.72x$ dollars. Say the single discount is of $k$. Then $(1-k)x=0.72x$. So $k=0.28$, or $k=28\%$. So select $\boxed{D}$.

 


	\item (\textit{AHSME 1950}) A man buys a house for $10000$ dollars and rents it. He puts $12\frac{1}{2}\%$ of each month's rent aside for repairs and upkeep; pays $325$ dollars a year taxes and realizes $5\frac{1}{2}\%$ on his investment. The monthly rent is:


 $\textbf{(A)}\ \$64.82 \qquad\textbf{(B)}\ \$83.33 \qquad\textbf{(C)}\ \$72.08\qquad\textbf{(D)}\ \$45.83 \qquad\textbf{(E)}\ \$177.08$

 \textbf{Answer: }B



\textbf{Solution}

$12\frac{1}{2}\%$ is the same as $\frac{1}{8}$, so the man sets one eighth of each month's rent aside, so he only gains $\frac{7}{8}$ of his rent. He also pays \$325 each year, and he realizes $5.5\%$, or \$550, on his investment. Therefore he must have collected a total of \$325 +\$550 = \$875 in rent. This was for the whole year, so he collected $\frac{875}{12}$ dollars each month as rent. This is only $\frac{7}{8}$ of the monthly rent, so the monthly rent in dollars is $\frac{875}{12}\cdot \frac{8}{7}=\boxed{\textbf{(B)}\ \$83.33}$.

 


	\item (\textit{AHSME 1950}) The equation $x+\sqrt{x-2}=4$ has:


 $\textbf{(A)}\ \text{2 real roots}\qquad\textbf{(B)}\ \text{1 real and 1 imaginary root}\qquad\textbf{(C)}\ \text{2 imaginary roots}\qquad\textbf{(D)}\ \text{No roots}\qquad\textbf{(E)}\ \text{1 real root}$

 \textbf{Answer: }E



\textbf{Solution 1}

$x + \sqrt{x-2} = 4$ Original Equation

 $\sqrt{x-2} = 4 - x$ Subtract x from both sides

 $x-2 = 16 - 8x + x^2$ Square both sides

 $x^2 - 9x + 18 = 0$ Get all terms on one side

 $(x-6)(x-3) = 0$ Factor

 $x = \{6, 3\}$

 If you put down A as your answer, it's wrong. You need to check for extraneous roots.

 $6 + \sqrt{6 - 2} = 6 + \sqrt{4} = 6 + 2 = 8 \ne 4$
$3 + \sqrt{3-2} = 3 + \sqrt{1} = 3 + 1 = 4 \checkmark$

 There is $\boxed{\textbf{(E)} \text{1 real root}}$




\textbf{Solution 2}

It's not hard to note that $x=3$ simply works, as $3 + \sqrt{1} = 4$. But, $x$ is increasing, and $\sqrt{x-2}$ is increasing, so $3$ is the only root. If $x < 3$, $x + \sqrt{x-2} < 4$, and similarly if $x > 3$, then $x + \sqrt{x-2} > 4$. Imaginary roots don't have to be considered because squaring the equation gives a quadratic, and quadratics either have only imaginary or only real roots.

 



\textbf{Solution 3}

We can create symmetry in the equation:
 \[x+\sqrt{x-2} = 4\]
 \[x-2+\sqrt{x-2} = 2.\]Let $y = \sqrt{x-2}$, then we have
 \[y^2+y-2 = 0\]
 \[(y+2)(y-1) = 0\]The two roots are $\sqrt{x-2} = -2, 1$.

 Notice, that the first root is extraneous as the range for the square root function is always the non-negative numbers (remember, negative numbers in square roots give imaginary numbers - imaginary numbers in square roots don't give negative numbers); thus, the only real root for $x$ occurs for the second root; squaring both sides and solving for $x$ gives $x=3 \Rightarrow \boxed{\textbf{(E)} \text{1 real root}}$.

 


	\item (\textit{AHSME 1950}) The value of $\log_{5}\frac{(125)(625)}{25}$ is equal to:


 $\textbf{(A)}\ 725\qquad\textbf{(B)}\ 6\qquad\textbf{(C)}\ 3125\qquad\textbf{(D)}\ 5\qquad\textbf{(E)}\ \text{None of these}$

 \textbf{Answer: }D



\textbf{Solution 1}

$\log_{5}\frac{(125)(625)}{25}$ can be simplified to $\log_{5}\ (125)(25)$ since $25^2 = 625$. $125 = 5^3$ and $5^2 = 25$ so $\log_{5}\ 5^5$ would be the simplest form. In $\log_{5}\ 5^5$, $5^x = 5^5$. Therefore, $x = 5$ and the answer is $\boxed{\mathrm{(D)}\ 5}$.

 



\textbf{Solution 2}

$\log_{5}\frac{(125)(625)}{25}$ can be also represented as $\log_{5}\frac{(5^3)(5^4)}{5^2}= \log_{5}\frac{(5^7)}{5^2}= \log_{5} 5^5$ which can be solved to get $\boxed{\mathrm{(D)}\ 5}$.

 


	\item (\textit{AHSME 1950}) if $\log_{10}{m}= b-\log_{10}{n}$, then $m=$
$\textbf{(A)}\ \frac{b}{n}\qquad\textbf{(B)}\ bn\qquad\textbf{(C)}\ 10^{b}n\qquad\textbf{(D)}\ b-10^{n}\qquad\textbf{(E)}\ \frac{10^{b}}{n}$

 \textbf{Answer: }E



\textbf{Solution 1}

We have $b=\log_{10}{10^b}$. Substituting, we find $\log_{10}{m}= \log_{10}{10^b}-\log_{10}{n}$. Using $\log{a}-\log{b}=\log{\dfrac{a}{b}}$, the left side becomes $\log_{10}{\dfrac{10^b}{n}}$. Because $\log_{10}{m}=\log_{10}{\dfrac{10^b}{n}}$, $m=\boxed{\mathrm{(E) }\dfrac{10^b}{n}}$.

 



\textbf{Solution 2}

adding $\log_{10} n$ to both sides:
 \[\log_{10} m + \log_{10} n=b\]using the logarithm property: $\log_a {b} + \log_a {c}=\log_a{bc}$
 \[\log_{10} {mn}=b\]rewriting in exponential notation:
 \[10^b=mn\]
 \[m=\boxed{\mathrm{(E) }\dfrac{10^b}{n}}\]~Vndom

 



\textbf{Solution 3}

More simply, we can just simulate the problem, if we have $m = 10$, that means the right side must be 1, so the only way we can achieve that with distinct $n$, is if $b = 3$, and $n = 100$. With this we can look through the different answer choices substituting in $b$, $n$, and $m$, and find that $\boxed{\mathrm{(E)}}$ is the only one that satisfies the question.

 


	\item (\textit{AHSME 1950}) A car travels $120$ miles from $A$ to $B$ at $30$ miles per hour but returns the same distance at $40$ miles per hour. The average speed for the round trip is closest to:


 $\textbf{(A)}\ 33\text{ mph}\qquad\textbf{(B)}\ 34\text{ mph}\qquad\textbf{(C)}\ 35\text{ mph}\qquad\textbf{(D)}\ 36\text{ mph}\qquad\textbf{(E)}\ 37\text{ mph}$

 \textbf{Answer: }B



\textbf{Solution}

The car takes $120 \text{ miles }\cdot\dfrac{1 \text{ hr }}{30 \text{ miles }}=4 \text{ hr}$ to get from $A$ to $B$. Also, it takes $120 \text{ miles }\cdot\dfrac{1 \text{ hr }}{40 \text{ miles }}=3 \text{ hr}$ to get from $B$ to $A$. Therefore, the average speed is $\dfrac{240\text{ miles }}{7 \text{ hr}}=34\dfrac{2}{7}\text{ mph}$, which is closest to $\boxed{\textbf{(B)}\ 34\text{ mph}}$.

 


	\item (\textit{AHSME 1950}) Two boys $A$ and $B$ start at the same time to ride from Port Jervis to Poughkeepsie, $60$ miles away. $A$ travels $4$ miles an hour slower than $B$.  $B$ reaches Poughkeepsie and at once turns back meeting $A$ $12$ miles from Poughkeepsie. The rate of $A$ was:


 $\textbf{(A)}\ 4\text{ mph}\qquad\textbf{(B)}\ 8\text{ mph}\qquad\textbf{(C)}\ 12\text{ mph}\qquad\textbf{(D)}\ 16\text{ mph}\qquad\textbf{(E)}\ 20\text{ mph}$

 \textbf{Answer: }B



\textbf{Solution}

Let the speed of boy $A$ be $a$, and the speed of boy $B$ be $b$. Notice that $A$ travels $4$ miles per hour slower than boy $B$, so we can replace $b$ with $a+4$.

 Now let us see the distances that the boys each travel. Boy $A$ travels $60-12=48$ miles, and boy $B$ travels $60+12=72$ miles. Now, we can use $d=rt$ to make an equation, where we set the time to be equal: 
 \[\frac{48}{a}=\frac{72}{a+4}\]Cross-multiplying gives $48a+192=72a$. Isolating the variable $a$, we get the equation $24a=192$, so $a=\boxed{\textbf{(B) }8 \text{ mph}}$.

 \textbf{Alternate Solution} Note that $A$ travels $60-12=48$ miles in the time it takes $B$ to travel $60+12=72$ miles. Thus, $B$ travels $72-48=24$ more miles in the given time, meaning $\frac{24\text{miles}}{4\text{miles}/\text{hour}} = 6 \text{hours}$ have passed, as $B$ goes $4$ miles per hour faster. Thus, $A$ travels $48$ miles per $6$ hours, or $8$ miles per hour.

 


	\item (\textit{AHSME 1950}) A manufacturer built a machine which will address $500$ envelopes in $8$ minutes. He wishes to build another machine so that when both are operating together they will address $500$ envelopes in $2$ minutes. The equation used to find how many minutes $x$ it would require the second machine to address $500$ envelopes alone is:


 $\textbf{(A)}\ 8-x=2\qquad\textbf{(B)}\ \frac{1}{8}+\frac{1}{x}=\frac{1}{2}\qquad\textbf{(C)}\ \frac{500}{8}+\frac{500}{x}=500\qquad\textbf{(D)}\ \frac{x}{2}+\frac{x}{8}=1\qquad\\ \textbf{(E)}\ \text{none of these answers}$

 \textbf{Answer: }B



\textbf{Solution 1}

We first represent the first machine's speed in per $2$ minutes: $125 \text{ envelopes in }2\text{ minutes}$. Now, we know that the speed per $2$ minutes of the second machine is 
 \[500-125=375 \text{ envelopes in }2\text{ minutes}\]Now we can set up a proportion to find out how many minutes it takes for the second machine to address $500$ papers:

 $\frac{375}{2}=\frac{500}{x}$. Solving for $x$, we get $x=\frac{8}{3}$ minutes.Now that we know the speed of the second machine, we can just plug it in each option to see if it equates. We see that $\boxed{\textbf{(B)}\ \dfrac{1}{8}+\dfrac{1}{x}=\dfrac{1}{2}}$ works.

 



\textbf{Solution 2}

First, notice that the number of envelopes addressed does not matter, because it stays constant throughout the problem. Next, notice that we are talking about combining two speeds, so we use the formula $\frac{1}{\frac{1}{a}+\frac{1}{b}}=c$, where $a,b$ are the respective times independently, and $c$ is the combined time. Plugging in for $a$ and $c$, we get $\frac{1}{\frac{1}{8}+\frac{1}{b}}=2$.

 Taking a reciprocal, we finally get $\boxed{\textbf{(B)}\ \dfrac{1}{8}+\dfrac{1}{x}=\dfrac{1}{2}}$.

 


	\item (\textit{AHSME 1950}) From a group of boys and girls, $15$ girls leave. There are then left two boys for each girl. After this $45$ boys leave. There are then $5$ girls for each boy. The number of girls in the beginning was:


 $\textbf{(A)}\ 40\qquad\textbf{(B)}\ 43\qquad\textbf{(C)}\ 29\qquad\textbf{(D)}\ 50\qquad\textbf{(E)}\ \text{None of these}$

 \textbf{Answer: }A



\textbf{Solution}

Let us represent the number of boys $b$, and the number of girls $g$.

 From the first sentence, we get that $2(g-15)=b$.

 From the second sentence, we get $5(b-45)=g-15$.

 Expanding both equations and simplifying, we get $2g-30 = b$ and $5b = g+210$.

 Substituting $b$ for $2g-30$, we get $5(2g-30)=g+210$. Solving for $g$, we get $g = \boxed{\textbf{(A)}\ 40}$.

 


	\item (\textit{AHSME 1950}) John ordered $4$ pairs of black socks and some additional pairs of blue socks. The price of the black socks per pair was twice that of the blue. When the order was filled, it was found that the number of pairs of the two colors had been interchanged. This increased the bill by $50\%$. The ratio of the number of pairs of black socks to the number of pairs of blue socks in the original order was:


 $\textbf{(A)}\ 4:1\qquad\textbf{(B)}\ 2:1\qquad\textbf{(C)}\ 1:4\qquad\textbf{(D)}\ 1:2\qquad\textbf{(E)}\ 1:8$

 \textbf{Answer: }C



\textbf{Solution}

Let the number of blue socks be represented as $b$. We are informed that the price of the black sock is twice the price of a blue sock; let us assume that the price of one pair of blue socks is $1$. That means the price of one pair of black socks is $2$.

 Now from the third and fourth sentence, we see that $1.5(2(4)+1(b))=1(4)+2(b)$. Simplifying gives $b=16$. This means the ratio of the number of pairs of black socks and the number of pairs of blue socks is $\boxed{\textbf{(C)}\ 1:4}$.

 


	\item (\textit{AHSME 1950}) A $25$ foot ladder is placed against a vertical wall of a building. The foot of the ladder is $7$ feet from the base of the building. If the top of the ladder slips $4$ feet, then the foot of the ladder will slide:


 $\textbf{(A)}\ 9\text{ ft}\qquad\textbf{(B)}\ 15\text{ ft}\qquad\textbf{(C)}\ 5\text{ ft}\qquad\textbf{(D)}\ 8\text{ ft}\qquad\textbf{(E)}\ 4\text{ ft}$

 \textbf{Answer: }D



\textbf{Solution 1}

By the Pythag theorem, triple $(7,24,25)$, the point where the ladder meets the wall is $24$ feet above the ground. When the ladder slides, it becomes $20$ feet above the ground. By the $(15,20,25)$ Pythagorean triple, The foot of the ladder is now $15$ feet from the building. Thus, it slides $15-7 = \boxed{\textbf{(D)}\ 8\text{ ft}}$.

 



\textbf{Solution 2}

We can observe that the above setup forms a right angled triangles whose base is 7ft and whose hypotenuse is 25ft taking the height to be x ft.

 
 \[x^2 + 7^2 = 25^2\]
 \[x^2 = 625 - 49\]
 \[x^2 = 576\]
 \[x = 24\]

 Since the top of the ladder slipped by 4 ft the new height is $24 - 4 = 20 ft$. The base of the ladder has moved so the new base is say $(7+y)$. The hypotenuse remains the same at 25ft. So,

 
 \[20^2 + (7+y)^2 = 25^2\]
 \[400 + 49 + y^2 + 14y = 625\]
 \[y^2 + 14y - 176 = 0\]
 \[y^2 + 22y - 8y - 176\]
 \[x(y+22) - 8(y+22)\]
 \[(y-8)(y+22)\]

 Disregarding the negative solution to equation the solution to the problem is $\boxed{\textbf{(D)}\ 8\text{ ft}}$.

 


	\item (\textit{AHSME 1950}) The number of circular pipes with an inside diameter of $1$ inch which will carry the same amount of water as a pipe with an inside diameter of $6$ inches is:


 $\textbf{(A)}\ 6\pi\qquad\textbf{(B)}\ 6\qquad\textbf{(C)}\ 12\qquad\textbf{(D)}\ 36\qquad\textbf{(E)}\ 36\pi$

 \textbf{Answer: }D



\textbf{Solution}

It must be assumed that the pipes have an equal height.

 We can represent the amount of water carried per unit time by cross sectional area.Cross sectional of Pipe with diameter $6 in$, 
 \[\pi r^2 = \pi \cdot 3^2 = 9\pi\]

 Cross sectional area of pipe with diameter $1 in$

 \[\pi r^2 = \pi \cdot 0.5 \cdot 0.5 = \frac{\pi}{4}\] 

 So number of 1 in pipes required is the number obtained by dividing their cross sectional areas

 
 \[\frac{9\pi}{\frac{\pi}{4}} = 36\]

 So the answer is $\boxed{\textbf{(D)}\ 36}$.

 



\textbf{Solution 2 (Solution 1 but easier)}

Knowing that $d^2 \alpha A$, where $d$ is the pipe diameter and $A$ is the cross-sectional area we simply get $6^2 = 36 = \boxed{\textbf{(D)}\ 36}$. 

 This works because the diameter of one of the other pipes is $1$, which is not affected by powers.

 


	\item (\textit{AHSME 1950}) When the circumference of a toy balloon is increased from $20$ inches to $25$ inches, the radius is increased by:


 $\textbf{(A)}\ 5\text{ in}\qquad\textbf{(B)}\ 2\frac{1}{2}\text{ in}\qquad\textbf{(C)}\ \frac{5}{\pi}\text{ in}\qquad\textbf{(D)}\ \frac{5}{2\pi}\text{ in}\qquad\textbf{(E)}\ \frac{\pi}{5}\text{ in}$

 \textbf{Answer: }D



\textbf{Solution 1}

When the circumference of a circle is increased by a percentage, the radius is also increased by the same percentage (or else the ratio of the circumference to the diameter wouldn't be $\pi$ anymore)We see that the circumference was increased by $25\%$. This means the radius was also increased by $25\%$. The radius of the original balloon is $\frac{20}{2\pi}=\frac{10}{\pi}$. With the $25\%$ increase, it becomes $\frac{12.5}{\pi}$. The increase is $\frac{12.5-10}{\pi}=\frac{2.5}{\pi}=\boxed{\textbf{(D)}\ \dfrac{5}{2\pi}\text{ in}}$.

 



\textbf{Solution 2}

Circumference of a circle is $2 \pi r$ so the radius is $\frac{circumference}{2 \pi}$

 So radius of first circle

 
 \[2 \pi r = 20\]
 \[r = \frac{10}{\pi}\]

 Radius of second circle
 \[2 \pi r = 25\]
 \[r = \frac{25}{2 \pi}\]

 The difference of these radii is 

 
 \[\frac{25}{2 \pi} - \frac{10}{\pi} = \frac{5}{2 \pi}\]

 So the answer is $\boxed{\textbf{(D)}\ \dfrac{5}{2\pi}\text{ in}}$.

 



\textbf{Solution 3}

Let the radius of the circle with the larger circumference be $r_2$ and the circle with the smaller circumference be $r_1$. Calculating the ratio of the two
 \[\frac{r_2}{r_1}=\frac{25}{20}=\frac{5}{4}\]
 \[4r_2=5r_1\]
 \[4(r_2-r_1)=r_1\]
 \[r_2-r_1=\frac{r_1}{4}=\frac{\frac{20}{2\pi}}{4}=\frac{10}{4\pi}=\boxed{\textbf{(D)}\ \dfrac{5}{2\pi}\text{ in}}\]



	\item (\textit{AHSME 1950}) In triangle $ABC$, $AC=24$ inches, $BC=10$ inches, $AB=26$ inches. The radius of the inscribed circle is:


 $\textbf{(A)}\ 26\text{ in}\qquad\textbf{(B)}\ 4\text{ in}\qquad\textbf{(C)}\ 13\text{ in}\qquad\textbf{(D)}\ 8\text{ in}\qquad\textbf{(E)}\ \text{None of these}$

 \textbf{Answer: }B



\textbf{Solution}

The inradius is equal to the area divided by semiperimeter. The area is $\frac{(10)(24)}{2} = 120$ because it's a right triangle, as it's side length satisfies the Pythagorean Theorem. The semiperimeter is $30$. Therefore the inradius is $\boxed{\textbf{(B)}\ 4}$.

 



\textbf{Solution 2}

Since this is a right triangle, we have
 \[\frac{a+b-c}{2}=\boxed{4}\]

 - kante314

 



\textbf{Solution 3}

We know that the formula of in radius is area of triangle by its semiperimeter , hence it is a right triangle => area is 1/2×24×10=120 and semiperimeter is 60/2=30 => 120/30=4HK🗿

 


	\item (\textit{AHSME 1950}) A merchant buys goods at $25\%$ off the list price. He desires to mark the goods so that he can give a discount of $20\%$ on the marked price and still clear a profit of $25\%$ on the selling price. What percent of the list price must he mark the goods?


 $\textbf{(A)}\ 125\%\qquad\textbf{(B)}\ 100\%\qquad\textbf{(C)}\ 120\%\qquad\textbf{(D)}\ 80\%\qquad\textbf{(E)}\ 75\%$

 \textbf{Answer: }A



\textbf{Solution 1}

Let $x$ be the list price. Then the merchant pays $0.75x$ for the goods.

 Let $m$ be the marked price.  Then if the merchant sells the goods for $(1-0.2)m = 0.8m$, his profit will be $0.25 \cdot 0.8m = 0.2m$.  Therefore, he paid $0.8m - 0.2m = 0.6m = 0.75x$, so $m = \frac{0.75x}{0.6} = 1.25x$, that is, the marked price is $\boxed{\textbf{(A)}\ 125\%}$ of the list price.

 



\textbf{Solution 2}

Without loss of generality, suppose the list price is $100$. Then the merchant pays $100\cdot .75=75$. 

 Let $s$ be the selling price.  Then the profit is $s - 75 = 0.25s$, so $s = 100$.  Let $x$ be the marked price.  Then $0.8x = s = 100$, so $x = 125$, or $\boxed{\textbf{(A)}\ 125\%}$ of the list price.

 


	\item (\textit{AHSME 1950}) If $y =\log_{a}{x}$, $a>1$, which of the following statements is incorrect?


 $\textbf{(A)}\ \text{If }x=1,y=0\qquad\\ \textbf{(B)}\ \text{If }x=a,y=1\qquad\\ \textbf{(C)}\ \text{If }x=-1,y\text{ is imaginary (complex)}\qquad\\ \textbf{(D)}\ \text{If }0<x<1,y\text{ is always less than 0 and decreases without limit as }x\text{ approaches zero}\qquad\\ \textbf{(E)}\ \text{Only some of the above statements are correct}$

 \textbf{Answer: }E



\textbf{Solution}

Let us first check 

 $\textbf{(A)}\ \text{If }x=1,y=0$. Rewriting in exponential form yields $a^0=1$, which is always true.

 $\textbf{(B)}\ \text{If }x=a,y=1$. Rewriting in exponential form yields $a^1=a$, which is always true.

 $\textbf{(C)}\ \text{If }x=-1,y\text{ is imaginary (complex)}$. If $x = -1$, then $\log_a -1 = \frac{\ln -1}{\ln a} = \frac{i\pi}{\ln a}$, so this is true.

 $\textbf{(D)}\ \text{If }0<x<1,y\text{ is always less than 0 and decreases without limit as }x\text{ approaches zero}$. Rewriting: $a^y=x$ such that $x<a$. Well, a power of $a$ can be less than $a$ only if $y<1$. And we observe, $y$ has no lower asymptote, because it is perfectly possible to have $y$ be $-100000000$; in fact, the lower $y$ gets, $x$ approaches $0$. This is also correct.

 Since all four of the previous statements are true, the only false statement is $\boxed{\textbf{(E)}\  \text{Only some of the above statements are correct}}$.

 


	\item (\textit{AHSME 1950}) If the expression $\begin{vmatrix}a & c\\ d & b\end{vmatrix}$ has the value $ab-cd$ for all values of $a, b, c$ and $d$, then the equation $\begin{vmatrix}2x & 1\\ x & x\end{vmatrix}= 3$:


 $\textbf{(A)}\ \text{Is satisfied for only 1 value of }x\qquad\\ \textbf{(B)}\ \text{Is satisified for 2 values of }x\qquad\\ \textbf{(C)}\ \text{Is satisified for no values of }x\qquad\\ \textbf{(D)}\ \text{Is satisfied for an infinite number of values of }x\qquad\\ \textbf{(E)}\ \text{none of these}$

 \textbf{Answer: }B



\textbf{Solution}

Note that $\begin{vmatrix}2x & 1\\ x & x\end{vmatrix}=2x \cdot x - 1 \cdot x = 2x^2-x$, so $2x^2 - x=3$. Subtracting $3$ from both sides yields $2x^2-x-3=0$, which factors to $(2x-3)(x+1)=0$. Thus, $x=\dfrac{3}{2}$ or $x = -1$, so the equation $\boxed{\textbf{(B)}\ \text{Is satisified for 2 values of }x}$.

 Note: Alternatively, one may note that the equation is quadratic with a positive discriminant, so it will be satisfied for exactly two values of $x$.

 


	\item (\textit{AHSME 1950}) Given the series $2+1+\frac{1}{2}+\frac{1}{4}+...$ and the following five statements:


 (1) the sum increases without limit


 (2) the sum decreases without limit


 (3) the difference between any term of the sequence and zero can be made less than any positive quantity no matter how small


 (4) the difference between the sum and 4 can be made less than any positive quantity no matter how small


 (5) the sum approaches a limit


 Of these statements, the correct ones are:


 $\textbf{(A)}\ \text{Only }3\text{ and }4\qquad\textbf{(B)}\ \text{Only }5\qquad\textbf{(C)}\ \text{Only }2\text{ and }4\qquad\textbf{(D)}\ \text{Only }2,3\text{ and }4\qquad\textbf{(E)}\ \text{Only }4\text{ and }5$

 \textbf{Answer: }E



\textbf{Solution}

This series is a geometric series with common ratio $\frac{1}{2}$. Using the well-known formula for the sum of an infinite geometric series, we obtain that this series has a value of $2\cdot \frac{1}{1-\frac{1}{2}}=4$. It immediately follows that statements 1 and 2 are false while statements 4 and 5 are true. 3 is false because we cannot make \textbf{any} term of the sequence approach 0, though we can choose some term that is less than any given positive quantity. The correct answer is therefore $\boxed{\textbf{(E)}\ \text{Only }4\text{ and }5}$.

 


	\item (\textit{AHSME 1950}) The limit of $\frac{x^{2}-1}{x-1}$ as $x$ approaches $1$ as a limit is:


 $\textbf{(A)}\ 0\qquad\textbf{(B)}\ \text{Indeterminate}\qquad\textbf{(C)}\ x-1\qquad\textbf{(D)}\ 2\qquad\textbf{(E)}\ 1$

 \textbf{Answer: }D



\textbf{Solution 1}

Since $\lim \limits_{x\to 1} (x^2-1) = 0$ and $\lim \limits_{x\to 1} (x-1) = 0$, by L'Hôpital's rule $\lim \limits_{x\to 1}\frac{x^2-1}{x-1} = \lim \limits_{x\to 1}\frac{2x}{1} = \boxed{\textbf{(D)}\ 2}$.

 



\textbf{Solution 2}

If $x \neq 1$, then $\frac {x^2-1}{x-1} = \frac{(x+1)(x-1)}{x-1} = x+1$. 

 Therefore, as $x$ approaches $1$, $\frac {x^2-1}{x-1} = x + 1$ approaches $1 + 1 = \boxed{\textbf{(D)}\ 2}$.

 



\textbf{Solution 3 (Formal)}

Let $f(x) = \frac{x^2-1}{x-1}$.  Then if $x \neq 1$, $f(x) = \frac {x^2-1}{x-1} = \frac{(x+1)(x-1)}{x-1} = x+1$. 

 Let $\epsilon > 0$ be any positive real number and let $\delta = \epsilon$.  Then if $0 < |x-1| < \delta$, $|x-1| = |x+1-2| = |f(x)-2| < \delta = \epsilon$, so $\lim \limits_{x\to 1}\frac{x^2-1}{x-1} = \boxed{\textbf{(D)}\ 2}$ by the definition of the limit.

 -j314andrews

 


	\item (\textit{AHSME 1950}) The least value of the function $ax^2+bx+c$ ($a>0$) is:


 $\textbf{(A)}\ -\frac{b}{a}\qquad\textbf{(B)}\ -\frac{b}{2a}\qquad\textbf{(C)}\ b^{2}-4ac\qquad\textbf{(D)}\ \frac{4ac-b^{2}}{4a}\qquad\textbf{(E)}\ \text{None of these}$

 \textbf{Answer: }D



\textbf{Solution}

The vertex of a parabola is at $x=-\frac{b}{2a}$ for $ax^2+bx+c$. Because $a>0$, the vertex is a minimum. Therefore $a(-\frac{b}{2a})^2+b(-\frac{b}{2a})+c=a(\frac{b^2}{4a^2})-\frac{b^2}{2a}+c=\frac{b^2}{4a}-\frac{2b^2}{4a}+c=-\frac{b^2}{4a}+c=\frac{4ac}{4a}-\frac {b^2}{4a}=\frac{4ac-b^2}{4a} \Rightarrow \mathrm{(D)}$.

 



\textbf{Solution 2}

Using calculus, we find that the derivative of the function is $2ax + b$, which has a critical point at $\frac{-b}{2a}$. The second derivative of the function is $2a$; since $a > 0$, this critical point is a minimum. As in Solution 1, plug this value into the function to obtain $\boxed{\textbf{(D)}\ \dfrac{4ac-b^2}{4a}}$.

 


	\item (\textit{AHSME 1950}) The equation $x^{x^{x^{.^{.^.}}}}=2$ is satisfied when $x$ is equal to:


 $\textbf{(A)}\ \infty\qquad\textbf{(B)}\ 2\qquad\textbf{(C)}\ \sqrt[4]{2}\qquad\textbf{(D)}\ \sqrt{2}\qquad\textbf{(E)}\ \text{None of these}$

 \textbf{Answer: }D



\textbf{Solution 1:}

Taking the log, we get $\log_x 2 = x^{x^{x^{.^{.^.}}}}=2$, and $\log_x 2 = 2$. Solving for x, we get $2=x^2$, and $\sqrt{2}=x \Rightarrow \mathrm{(D)}$




\textbf{Solution 2:}

$x^{x^{x^{.^{.^.}}}}=2$ is the original equation. If we let $y=x^{x^{x^{.^{.^.}}}}$, then the equation can be written as $y=2$. This also means that $x^y=2$, considering that adding one $x$ to the start and then taking that $x$ to the power of $y$ does not have an effect on the equation, since $y$ is infinitely long in terms of $x$ raised to itself forever. It is already known that $y=2$ from what we first started with, so this shows that $x^y=x^2=2$. If $x^2=2$, then that means that $\sqrt{2}=x \Rightarrow \mathrm{(D)}$.

 This is just a faster way once you get used to it, instead of taking a log of the function.

 



\textbf{Solution 3:}

Given $x^{x^{x^{.^{.^.}}}}=2$, we see that $x$ is also raised to $x^{x^{x^{.^{.^.}}}}$, so $x^2 = 2$ $\Rightarrow x=\sqrt{2} \mathrm{{D}}$



	\item (\textit{AHSME 1950}) The sum to infinity of $\frac{1}{7}+\frac{2}{7^{2}}+\frac{1}{7^{3}}+\frac{2}{7^{4}}+...$ is:


 $\textbf{(A)}\ \frac{1}{5}\qquad\textbf{(B)}\ \frac{1}{24}\qquad\textbf{(C)}\ \frac{5}{48}\qquad\textbf{(D)}\ \frac{1}{16}\qquad\textbf{(E)}\ \text{None of these}$

 \textbf{Answer: }E



\textbf{Solution}

Note that this is $\frac{1}{7}(1+\frac{1}{49}+\frac{1}{49^2}+...)+\frac{2}{49}(1+\frac{1}{49}+...)=\frac{9}{49}(1+\frac{1}{49}+...)$. Using the formula for a geometric series, we find that this is $\frac{9}{49}(\frac{1}{1-\frac{1}{49}})=\frac{9}{49}(\frac{1}{\frac{48}{49}})=\frac{9}{49}(\frac{49}{48})=\frac{9}{48}=\frac{3}{16} \Rightarrow \mathrm{(E)}$



	\item (\textit{AHSME 1950}) The graph of $y=\log x$
$\textbf{(A)}\ \text{Cuts the }y\text{-axis}\qquad\\ \textbf{(B)}\ \text{Cuts all lines perpendicular to the }x\text{-axis}\qquad\\ \textbf{(C)}\ \text{Cuts the }x\text{-axis}\qquad\\ \textbf{(D)}\ \text{Cuts neither axis}\qquad\\ \textbf{(E)}\ \text{Cuts all circles whose center is at the origin}$

 \textbf{Answer: }C



\textbf{Solution}

The graph of $y = \log x$ includes $(1, 0)$, so it $\boxed{\textbf{(C)}\ \text{Cuts the }x\text{-axis}}$.

 Choice (A) is incorrect since $\log 0$ is undefined.

 Choice (B) is incorrect since $\log 0$ is undefined and $x = 0$ is perpendicular to the $x$-axis.

 Choice (D) is incorrect since $y = \log x$ intersects the $x$-axis.

 Choice (E) is incorrect since $y = \log x$ does not intersect a circle of radius $0.1$ centered at the origin.

 


	\item (\textit{AHSME 1950}) The number of diagonals that can be drawn in a polygon of $100$ sides is:


 $\textbf{(A)}\ 4850\qquad\textbf{(B)}\ 4950\qquad\textbf{(C)}\ 9900\qquad\textbf{(D)}\ 98\qquad\textbf{(E)}\ 8800$

 \textbf{Answer: }A



\textbf{Solution 1}

The $100$-gon has $\binom{100}{2}=4950$ distinct pairs of vertices. Of these pairs, $100$ determine sides of the polygon.  All other pairs determine diagonals, so the $100$-gon has $4950-100=\boxed{\textbf{(A)}\ 4850 }$ diagonals.

 



\textbf{Solution 2}

For the first endpoint of the diagonal, any of the $100$ vertices of the $100$-gon may be chosen.  For the second endpoint, any vertex that is neither the same as nor adjacent to the first endpoint may be chosen, for a total of $100 - 3 = 97$ choices.  However, this counts each diagonal twice as the order in which the endpoints are selected is irrelevant.  Therefore, the $100$-gon has $\frac{100 \cdot 97}{2} = \boxed{\textbf{(A)}\ 4850}$ diagonals.

 



\textbf{Solution 3}

It is well-known that an $n$-gon has $\frac{n(n-3)}{2}$ diagonals. Therefore, the $100$-gon has $\frac{100\cdot 97}{2} = \boxed{\textbf{(A)}\ 4850 }$ diagonals.

 -\href{https://artofproblemsolving.com/wiki/index.php/User:Cxsmi}{cxsmi}, edited by j314andrews

 


	\item (\textit{AHSME 1950}) In triangle $ABC$, $AB=12$, $AC=7$, and $BC=10$. If sides $AB$ and $AC$ are doubled while $BC$ remains the same, then:


 $\textbf{(A)}\ \text{The area is doubled}\qquad\\ \textbf{(B)}\ \text{The altitude is doubled}\qquad\\ \textbf{(C)}\ \text{The area is four times the original area}\qquad\\ \textbf{(D)}\ \text{The median is unchanged}\qquad\\ \textbf{(E)}\ \text{The area of the triangle is 0}$

 \textbf{Answer: }E



\textbf{Solution}

After doubling, $AB = 24$ and $AC= 14$.  Notice that now $AB = AC + BC$. Therefore, $C$ lies on $\overline{AB}$ and $\boxed{\textbf{(E)}\ \text{The area of the triangle is 0}}$.

 


	\item (\textit{AHSME 1950}) A rectangle inscribed in a triangle has its base coinciding with the base $b$ of the triangle. If the altitude of the triangle is $h$, and the altitude $x$ of the rectangle is half the base of the rectangle, then:


 $\textbf{(A)}\ x=\frac{1}{2}h\qquad\textbf{(B)}\ x=\frac{bh}{b+h}\qquad\textbf{(C)}\ x=\frac{bh}{2h+b}\qquad\textbf{(D)}\ x=\sqrt{\frac{hb}{2}}\qquad\\ \textbf{(E)}\ x=\frac{1}{2}b$

 \textbf{Answer: }C



\textbf{Solution 1:Similarity}

Draw the triangle.The small triangle formed by taking away the rectangle and the two small portions left is similar to the big triangle.Hence, the proportions of the heights is equal to the proportions of the sides. 

 In particular, we get
 \[\dfrac{2x}{b} = \dfrac{h - x}{h} \implies 2xh = bh - bx \implies (2h + b)x = bh \implies x = \dfrac{bh}{2h + b}\]The answer is $\boxed{\textbf{(C)}}$.

 



\textbf{Solution 2: Area}

The smaller triangle at the top has an area of 
 \[\dfrac{2x(h-x)}{2}=x(h-x)\]

 Because the two other pieces are complementary triangles, we can add them together to create one triangle with area
 \[\dfrac{x(b-2x)}{2}\]

 Lastly the area of the rectangle is $2x^2$. These areas together sum to the area of the big rectangle, which is $\dfrac{bh}{2}$. Hence,
 \[\dfrac{bh}{2}=2x^2 + \dfrac{x(b-2x)}{2}+x(h-x)\]

 Solving we get that 
 \[x = \dfrac{bh}{2h + b}\]The answer is $\boxed{\textbf{(C)}}$.

 


	\item (\textit{AHSME 1950}) A point is selected at random inside an equilateral triangle. From this point perpendiculars are dropped to each side. The sum of these perpendiculars is:


 $\textbf{(A)}\ \text{Least when the point is the center of gravity of the triangle}\qquad\\ \textbf{(B)}\ \text{Greater than the altitude of the triangle}\qquad\\ \textbf{(C)}\ \text{Equal to the altitude of the triangle}\qquad\\ \textbf{(D)}\ \text{One-half the sum of the sides of the triangle}\qquad\\ \textbf{(E)}\ \text{Greatest when the point is the center of gravity}$

 \textbf{Answer: }C



\textbf{Solution}

Let $A$, $B$, and $C$ be the vertices of the equilateral triangle and $P$ be the randomly selected point inside $\triangle ABC$.  Let $A'$, $B'$, $C'$ be the feet of the perpendiculars from $P$ to $\overline{BC}$, $\overline{AC}$, and $\overline{AB}$ respectively.  

 Let $s$ and $h$ be the side length and altitude of $\triangle ABC$ respectively.  Let $h_1 = PA'$, $h_2 = PB'$, and $h_3 = PC'$.  

 Then $[ABC] = \frac{1}{2}sh$.  Also, $[ABC] = [PAB] + [PAC] + [PBC] = \frac{1}{2}sh_1 + \frac{1}{2}sh_2 + \frac{1}{2}sh_3 = \frac{1}{2}s(h_1+h_2+h_3)$.

 Therefore, $[ABC] = \frac{1}{2}sh = \frac{1}{2}s(h_1+h_2+h_3)$, so $h = h_1+h_2+h_3$, that is, the sum of the perpendiculars from $P$ is $\boxed{\textbf{(C)}\ \textrm{Equal to the altitude of the triangle}}$.

 \textbf{Note} This result is also known as Viviani's theorem.

 


	\item (\textit{AHSME 1950}) A triangle has a fixed base $AB$ that is $2$ inches long. The median from $A$ to side $BC$ is $1 \frac{1}{2}$ inches long and can have any position emanating from $A$. The locus of the vertex $C$ of the triangle is:


 $\textbf{(A)}\ \text{A straight line }AB,1\frac{1}{2}\text{ inches from }A\qquad\\ \textbf{(B)}\ \text{A circle with }A\text{ as center and radius }2\text{ inches}\qquad\\ \textbf{(C)}\ \text{A circle with }A\text{ as center and radius }3\text{ inches}\qquad\\ \textbf{(D)}\ \text{A circle with radius }3\text{ inches and center }4\text{ inches from }B\text{ along }BA\qquad\\ \textbf{(E)}\ \text{An ellipse with }A\text{ as focus}$

 \textbf{Answer: }D



\textbf{Solution 1}

The locus of the median's endpoint on $BC$ is the circle about $A$ and of radius $1\frac{1}{2}$ inches.  The locus of the vertex $C$ is then the circle twice as big and twice as far from $B$, i.e. of radius $3$ inches and with center $4$ inches from $B$ along $BA$ which means that our answer is: $\textbf{(D)}$.

 


 



\textbf{Solution 2}

Let $A(a,y_A)$, $B(b,y_B)$ and $C(x,y)$.

 Hence, $D\left(\frac{x+b}{2},\frac{y+y_B}{2}\right)$ is a midpoint of $BC$.

 Thus, the equation of needed locus is
 \[\left(\frac{x+b}{2}-a\right)^2+\left(\frac{y+y_B}{2}-y_A\right)^2=\left(\frac{3}{2}\right)^2,\]which is equation of the circle:
 \[(x-(2a-b))^2+(y-(2y_A-y_B))^2=3^2.\]

 Thus, D) is valid because 
 \[\sqrt{(2a-b-b)^2+(2y_A-y_B-y_B)^2}=2\sqrt{(a-b)^2+(y_A-y_b)^2}=2AB=4.\]




\textbf{Solution 3 (Complex Numbers)}

Place $\triangle ABC$ in the complex plane such that $A = 0$ and $B = 2$.  Let $M$ the midpoint of $BC$.  Since $AM = 1.5$, the locus of $M$ is $1.5z$ over all $z$ such that $|z| = 1$.  Since $M = \frac{B+C}{2}$, $C = 2M-B = 3z-2$.  Therefore, the locus of $C$ is a circle of radius $3$ centered at $-2$, which is $4$ units from $B$ along $BA$, thus the answer is $\boxed{(\textbf{D})}$.

 -j314andrews

 


 


	\item (\textit{AHSME 1950}) A privateer discovers a merchantman $10$ miles to leeward at $11\text{:}45 \text{ a.m.}$ and with a good breeze bears down upon her at $11$ mph, while the merchantman can only make $8$ mph in his attempt to escape. After a two hour chase, the top sail of the privateer is carried away; she can now make only $17$ miles while the merchantman makes $15$. The privateer will overtake the merchantman at:


 $\textbf{(A)}\ 3\text{:}45\text{ p.m.}\qquad\textbf{(B)}\ 3\text{:}30\text{ p.m.}\qquad\textbf{(C)}\ 5\text{:}00\text{ p.m.}\qquad\textbf{(D)}\ 2\text{:}45\text{ p.m.}\qquad\textbf{(E)}\ 5\text{:}30\text{ p.m.}$

 \textbf{Answer: }E



\textbf{Solution 1}

Assume that the two boats are traveling along the positive real number line, with the merchantman starting at the number $10$ and the privateer starting at the number $0$. After two hours the merchantman is at $26$ while the privateer is at $22$. When the top sail of the privateer is carried away, the speed of the merchantman is unaffected; he still travels at $8$ miles per hour. Therefore the privateer travels at $17\cdot \frac{8}{15}=\frac{136}{15}=9\frac{1}{15}$ miles per hour. The remaining time $t$ in hours it takes for the privateer to catch up to the merchantman satisfied the equation

 
 \[26+8t=22+\frac{136}{15}t\]

 Simplification yields the equation $\frac{16}{15}t=4$, which shows that $t=\frac{15}{4}=3\frac{3}{4}$. It therefore takes a total of $5\frac{3}{4}$ hours for the privateer to catch the merchantman, so this will happen at $\boxed{\textbf{(E)}\ 5\text{:}30\text{ p.m.}}$




\textbf{Solution 2}

At first the privateer travels $3$ mph faster than the merchantman, so over the first $2$ hours, the privateer gains $2 \cdot 3 = 6$ miles on the merchantman.  Therefore, after $2$ hours, the privateer is $10 - 6 = 4$ miles behind the merchantman.

 After losing her top sail, the privateer travels at $\frac{17}{15}$ times the speed of the merchantman, or $8 \cdot \frac{17}{15} = \frac{136}{15}$ mph.  Since the privateer now travels $\frac{136}{15} - 8 = \frac{16}{15}$ mph faster than the merchantman, it will take $4 \div \frac{16}{15} = \frac{15}{4}$ more hours for the privateer to catch up.  

 Therefore, it will take a total of $2 + \frac{15}{4} = \frac{23}{4}$ hours, or $5$ hours and $45$ minutes for the privateer to reach the merchantman, so it will be $\boxed{\textbf{(E)}\ 5\text{:}30\text{ p.m.}}$

 -j314andrews

 


\end{enumerate}